% Figura (Problemas 1.2 y 1.9): superficies gaussianas y gráficas E(r), phi(r) para esfera maciza y cáscara.
\begin{figure}[ht]
\centering
\begin{minipage}[t]{0.30\textwidth}\centering
\begin{tikzpicture}[baseline=(current bounding box.north), scale=0.95]
  \shade[ball color=azulclaro!50, opacity=0.30] (0,0) circle (1.2);
  \draw[azulprof, line width=0.8pt] (0,0) circle (1.2);
  \draw[verdeteo, densely dashed, line width=0.7pt] (0,0) circle (0.7);
  \draw[verdeteo, densely dashed, line width=0.7pt] (0,0) circle (2.0);
  \foreach \a in {0,45,...,315}{\draw[-{Stealth[length=4pt]}, cm1, line width=0.6pt] (\a:2.0) -- (\a:2.45);}
  \foreach \a in {0,45,...,315}{\draw[-{Stealth[length=3pt]}, cm1, line width=0.6pt] (\a:0.7) -- (\a:0.95);}
  \foreach \p in {(-0.4,0.9),(0.9,0.3),(-0.95,-0.35),(0.3,-0.95)}{\node[font=\tiny, text=azulacero] at \p {$+$};}
  \draw[grispizarra, ->, line width=0.4pt] (0,0) -- (-125:1.2) node[pos=0.5, left, font=\scriptsize, text=azulprof] {$a$};
  \node[font=\scriptsize, text=verdeteo, anchor=south west] at (40:0.62) {$r$};
  \node[font=\scriptsize, text=verdeteo, anchor=north west] at (-45:2.05) {$r$};
  \node[font=\scriptsize, text=verdeteo] at (0,2.7) {$S_{\mathrm{G}}$};
\end{tikzpicture}\\[2pt]{\small\scshape (a)}
\end{minipage}\hfill
\begin{minipage}[t]{0.33\textwidth}\centering
\begin{tikzpicture}[baseline=(current bounding box.north)]
  \begin{axis}[tema, width=\textwidth, height=4.2cm, xlabel={$r/a$}, ylabel={$E/E(a)$}, xmin=0, xmax=3, ymin=0, ymax=1.15,
     xtick={0,1,2,3}, ytick={0,0.5,1}, domain=0:3, samples=120]
    \addplot[azulprof, line width=1.1pt, domain=0:1] {x};
    \addplot[azulprof, line width=1.1pt, domain=1:3] {1/x^2};
    \addplot[grispizarra, densely dashed] coordinates {(1,0) (1,1)};
    \node[font=\scriptsize, text=azulprof, anchor=south west] at (axis cs:0.15,0.35) {$\propto r$};
    \node[font=\scriptsize, text=azulprof, anchor=south west] at (axis cs:1.3,0.55) {$\propto 1/r^{2}$};
  \end{axis}
\end{tikzpicture}\\[2pt]{\small\scshape (b)}
\end{minipage}\hfill
\begin{minipage}[t]{0.33\textwidth}\centering
\begin{tikzpicture}[baseline=(current bounding box.north)]
  \begin{axis}[tema, width=\textwidth, height=4.2cm, xlabel={$r/a$}, ylabel={$\phi/\phi(a)$}, xmin=0, xmax=3, ymin=0, ymax=1.6,
     xtick={0,1,2,3}, ytick={0,0.5,1,1.5}, domain=0:3, samples=120]
    \addplot[azulacero, line width=1.1pt, domain=0:1] {(3-x^2)/2};
    \addplot[azulacero, line width=1.1pt, domain=1:3] {1/x};
    \addplot[grispizarra, densely dashed] coordinates {(1,0) (1,1)};
    \node[font=\scriptsize, text=azulacero, anchor=south west] at (axis cs:1.35,0.5) {$\propto 1/r$};
    \node[font=\scriptsize, text=azulacero, anchor=south west] at (axis cs:0.05,0.75) {$\tfrac{3}{2}-\tfrac{r^{2}}{2a^{2}}$};
  \end{axis}
\end{tikzpicture}\\[2pt]{\small\scshape (c)}
\end{minipage}
\caption{Problema 1.2. (a)~Corte de la esfera uniformemente cargada y de las dos superficies gaussianas esféricas (a trazos), una
interior ($r<a$) y otra exterior ($r>a$). (b)~El campo crece linealmente dentro y decae como $1/r^{2}$ fuera, con
valor máximo $E(a)$ en la superficie. (c)~El potencial, máximo en el centro, es parabólico dentro y decae como $1/r$ fuera; $E$ y $\phi$ son continuos en $r=a$.}
\label{fig:sem-gauss-esfera}
\end{figure}
